% source: David Mumford, "Abelian Varieties", TIFR Studies in Mathematics 5, Oxford UP, 1970
% pdf_page: 0063
% doc_page: 52 (Chapter II, Section 5, "Cohomology and base change")
% retrieved: 2026-07-11
% transcribed_by: claude-fable-5 (vision, rendered at 200dpi; scanned book, no text layer)

% requires: amsmath (bmatrix, \text, \xrightarrow), amssymb (\mathfrak), mathrsfs (\mathscr)

\DeclareMathOperator{\Ima}{Im}
\DeclareMathOperator{\Ker}{Ker}
\DeclareMathOperator{\Spec}{Spec}

% [running head: ABELIAN VARIETIES, page 52]

% [continuation of the statement of Lemma 2 from the previous page]
\textit{be a homomorphism of coherent locally free $\mathcal{O}_Y$-sheaves.
If $\dim_{k(y)}[\Ima(\phi \otimes k(y))]$ is locally constant, then there are
splittings:}
\[
  \mathscr{F} \simeq \mathscr{F}_1 \oplus \mathscr{F}_2
\]
\[
  \mathfrak{D} \simeq \mathfrak{D}_1 \oplus \mathfrak{D}_2
\]
\textit{such that $\phi|_{\mathscr{F}_1} = (0)$,
$\Ima(\phi) \subset \mathfrak{D}_1$, and
$\phi : \mathscr{F}_2 \to \mathfrak{D}_1$ is an isomorphism, i.e.}
\[
  \phi =
  \begin{bmatrix}
    0 & \text{isom.} \\
    0 & 0
  \end{bmatrix}
\]

\paragraph{Proof.}
By Lemma 1, $\mathfrak{D}/\phi(\mathscr{F})$ is locally free. If
$Y = \Spec(A)$, $M = \Gamma(Y, \mathscr{F})$, $N = \Gamma(Y, \mathfrak{D})$,
then this means that $N/\phi(M)$ is $A$-projective. Therefore $N$ splits into
the direct sum of $\phi(M)$ and a second submodule isomorphic to
$N/\phi(M)$. Or, in sheaves, $\mathfrak{D} \simeq \mathfrak{D}_1 \oplus \mathfrak{D}_2$,
where $\mathfrak{D}_1 = \Ima(\phi)$. Moreover, this shows that $\phi(M)$ is
$A$-projective, too, so $M$ splits into the direct sum of $\Ker(\phi)$ and a
second submodule isomorphic to $\phi(M)$. Or, in sheaves,
$\mathscr{F} \simeq \mathscr{F}_1 \oplus \mathscr{F}_2$, where
$\phi(\mathscr{F}_1) = (0)$,
$\phi : \mathscr{F}_2 \xrightarrow{\ \sim\ } \mathfrak{D}_1$.

Now assume (i) holds. Let $K^{\cdot}$ be the complex given by the theorem.
As in the proof of Corollary 1, $\dim[\Ima(d^{p-1} \otimes k(y))]$
% [in the print the first "k(y)" is smudged and reads like "h(y)"; clearly k(y) as in the parallel expression on the next line]
and $\dim[\Ima(d^p \otimes k(y))]$ are locally constant. By Lemma 2, applied
first to $d_p : K^p \to K^{p+1}$, and second to
$d_{p-1} : K^{p-1} \to \Ker(d_p)$, we get splittings into projective modules:
\[
  \begin{array}{ccccc}
    Z_{p-1} \oplus K'_{p-1} & & B_p \oplus H_p \oplus K'_p & & B_{p+1} \oplus K'_{p+1} \\
    \| & & \| & & \| \\
    K_{p-1} & \longrightarrow & K_p & \longrightarrow & K_{p+1}
  \end{array}
\]
where $Z_{p-1} = \Ker(d_{p-1})$, $d_{p-1} : K'_{p-1} \to B_p$ is an
isomorphism, $B_p \oplus H_p = \Ker(d_p)$, and $d_p : K'_p \to B_{p+1}$ is an
isomorphism. It follows immediately that
\[
  H^p(K^{\cdot} \otimes_A B) \simeq H_p \otimes_A B \simeq
  H^p(K^{\cdot}) \otimes_A B, \text{ all } B
\]
and
$H^{p-1}(K^{\cdot} \otimes_A B) \simeq
 Z_{p-1} \otimes_A B / \Ima(d_{p-2} \otimes B) \simeq
 H^{p-1}(K^{\cdot}) \otimes_A B$,
all $B$. This proves (ii).
